\documentclass[12pt,a4paper]{article}

% ── Fonts and encoding (LuaLaTeX) ─────────────────────────────────────────────
\usepackage{fontspec}
\usepackage{unicode-math}
\setmainfont{Linux Libertine O}
\setmathfont{TeX Gyre Pagella Math}
\setmonofont[Scale=0.85]{Everson Mono}

% ── Layout ────────────────────────────────────────────────────────────────────
\usepackage[top=1in, bottom=1in, left=1.1in, right=1.1in, headheight=15pt]{geometry}
\usepackage{microtype}
\usepackage{parskip}
\usepackage[english]{babel}
\setlength{\emergencystretch}{3em}

% ── Headers ───────────────────────────────────────────────────────────────────
\usepackage{fancyhdr}
\pagestyle{fancy}
\fancyhf{}
\fancyhead[L]{\leftmark}
\fancyhead[R]{\thepage}
\renewcommand{\headrulewidth}{0.4pt}

% ── Section styling ───────────────────────────────────────────────────────────
\usepackage{titlesec}
\titleformat{\section}{\Large\bfseries}{\thesection}{1em}{}
\titleformat{\subsection}{\large\bfseries}{\thesubsection}{1em}{}
\titleformat{\subsubsection}{\normalsize\bfseries}{\thesubsubsection}{1em}{}

% ── Lists ─────────────────────────────────────────────────────────────────────
\usepackage{enumitem}
\setlist[itemize]{topsep=0.5ex, itemsep=0.25ex, parsep=0pt}
\setlist[enumerate]{topsep=0.5ex, itemsep=0.25ex, parsep=0pt}

% ── Math and tables ───────────────────────────────────────────────────────────
\usepackage{amsmath,amsthm}
\usepackage{mathtools}
\usepackage{booktabs}
\usepackage{array}
\usepackage{tabularx}
\usepackage[table]{xcolor}
\renewcommand{\arraystretch}{1.3}

% ── Code listings ─────────────────────────────────────────────────────────────
\usepackage{listings}
\definecolor{codebg}{RGB}{248,248,248}
\definecolor{codeframe}{RGB}{200,200,200}
\definecolor{codekeyword}{RGB}{0,0,180}
\definecolor{codecomment}{RGB}{0,128,0}
\lstset{
  basicstyle=\ttfamily\footnotesize,
  backgroundcolor=\color{codebg},
  frame=single,
  rulecolor=\color{codeframe},
  keywordstyle=\color{codekeyword}\bfseries,
  commentstyle=\color{codecomment}\itshape,
  breaklines=true,
  captionpos=b,
  morekeywords={axiom,def,theorem,noncomputable,structure,where,namespace,
                open,import,section,end,fun,exact,intro,simp,rfl,decide,
                constructor,obtain,refine,cases,by,have,exact,use,rw},
}

% ── Cross-references ──────────────────────────────────────────────────────────
\usepackage{hyperref}
\usepackage{cleveref}
\hypersetup{colorlinks=true, linkcolor=blue, urlcolor=cyan, citecolor=blue,
  pdfkeywords={SIC-POVM, Belnap multilattice, Stark conjecture, Hilbert Twelfth Problem, Zauner conjecture, Lean 4, paraconsistent logic, Weyl-Heisenberg group, ray class fields, machine-checked proof}}

% ── Theorem environments ──────────────────────────────────────────────────────
\newtheorem{theorem}{Theorem}[section]
\newtheorem{lemma}[theorem]{Lemma}
\newtheorem{proposition}[theorem]{Proposition}
\newtheorem{corollary}[theorem]{Corollary}
\newtheorem{definition}[theorem]{Definition}
\newtheorem{axiom_env}[theorem]{Axiom}
\newtheorem{remark}[theorem]{Remark}

% ── Custom notation ───────────────────────────────────────────────────────────
\newcommand{\BFour}{\mathbf{B}_4}
\newcommand{\WH}{\mathrm{WH}}
\newcommand{\SIC}{\mathrm{SIC}}
\newcommand{\Kd}{K_d}
\newcommand{\Fd}{F_d}
\newcommand{\md}{m_d}
\newcommand{\frob}{\mu \circ \delta = \mathrm{id}}
\newcommand{\Lean}{\textsc{Lean\,4}}

% ─────────────────────────────────────────────────────────────────────────────

\begin{document}

\title{\textbf{SIC-POVMs, a Stark Conjecture, and the 12th:\\
A Formalization via Paraconsistent Belnap Multilattices}}

\author{C.\ Lando\ Mills%}
\thanks{The author thanks Harry T.\ Larson; see Acknowledgements and \cite{Larson1961}.}}

\date{7 July 2026}

\maketitle

\begin{abstract}
This paper establishes, inside the \Lean{} proof assistant, a three-level formal identification.  
The levels are: \emph{(i)}~Belnap multilattice axioms for Weyl--Heisenberg covariant SIC-POVMs at $d=2^n$;  
\emph{(ii)}~the Zauner conjecture; and  
\emph{(iii)}~the mixed-signature Stark conjecture for the ray class field
$\Kd=\mathbb{Q}(\sqrt{(d-3)(d+1)})$, a real-quadratic case of Hilbert's Twelfth Problem.  
Fiducials are unit-normalized and satisfy $(d+1)|\langle\psi,D_{a,b}\psi\rangle|^2=1$.

The equivalence \texttt{hilbert\_embedding\_equiv\_zauner} is proved by \texttt{rfl}:
the Belnap embedding into $\mathbb{C}^{2^n}$ and the Zauner conjecture at $d=2^n$
are definitionally the same proposition.  
The Belnap skeleton (orbit size $4^n$, Frobenius closure $\frob$, join-equiangularity, Born rule)
contains zero sorries.  
Open arithmetic content is marked by named gap axioms for Stark units on WH frames; a proof of Stark
would close all three levels at once.

For dimension $d=12$ we prove \texttt{SICPOVM\_Exists 12} outright.  
We construct an exact fiducial in a finitely presented $\mathbb{Q}$-algebra, verify $143$ overlap identities
with \texttt{native\_decide}, and transfer everything to $\mathbb{C}^{12}$ along a ring homomorphism.  
The theorem \texttt{crystal\_forces\_d12\_sic} depends on no axiom beyond \Lean{}'s standard foundations
and compiler trust. This is, to our knowledge, the first machine-checked SIC-POVM existence in any dimension.

For the frontier dimension $d=2048=2^{11}$ the transport apparatus is formalized and sorry-free.  
It includes a forward map $\varphi\colon B^{\oplus 11}\to\mathbb{C}^{2048}$, a reduction $\psi$ with
$\psi\circ\varphi=\mathrm{id}$, a conditional reduction to Stark, and a non-real character obstruction
that blocks the false branch.  
Unconditional existence remains open; the machinery that surrounds it is closed.
\end{abstract}

\noindent\textbf{Keywords:} SIC-POVM \(\cdot\) Belnap multilattice \(\cdot\)
Stark conjecture \(\cdot\) Hilbert's Twelfth Problem \(\cdot\) Zauner conjecture
\(\cdot\) Lean 4 \(\cdot\) paraconsistent logic \(\cdot\) Weyl-Heisenberg group
\(\cdot\) ray class fields \(\cdot\) machine-checked proof

\newpage
\tableofcontents
\newpage

%──────────────────────────────────────────────────────────────────────────────
\section{Introduction}
\label{sec:intro}
%──────────────────────────────────────────────────────────────────────────────

Symmetric informationally complete positive operator-valued measures (SIC-POVMs)
are the minimal tomographically complete quantum measurements:
$d^2$ rank-1 projectors in $\mathbb{C}^d$ with constant pairwise Hilbert--Schmidt inner products.
Zauner conjectured in 1999 that such structures exist in every finite dimension \cite{Zauner1999}.
Numerical checks have pushed the frontier past $d \approx 2400$, but a general proof remains absent.

Appleby, Flammia, McConnell, Yard, and collaborators \cite{Appleby2013,Appleby2017,AFMY2020}
discovered that the fiducial vector coordinates inhabit specific ray class fields of real
quadratic fields.
For a given dimension $d$ the base field is $F_d = \mathbb{Q}(\sqrt{(d-3)(d+1)})$, and the coordinates
generate the ray class field $\Kd$ of $F_d$.  
This ties SIC-POVMs directly to Hilbert's Twelfth Problem
(constructing abelian extensions of number fields from special values of transcendental functions)
in the still-open real quadratic case.

We formalize these connections in \Lean{}\cite{Lean4} (v4.28.0, paraconsistent kernel fork \cite{p4rakernel})
using five modules:

\begin{itemize}
  \item \texttt{Imscribing.Millennium.SIC\_POVM\_Stark}: Weyl--Heisenberg operators in $\mathbb{C}^d$,
    the SIC-POVM definition, Stark arithmetic, and the conditional existence theorem (\S\ref{sec:stark}).
  \item \texttt{Imscribing.Paraconsistent.Shor.BelnapWHMultilattice}: the Belnap $\BFour$ multilattice,
    the product-lattice orbit theorem ($2^n$ proved), and the multilattice axioms for $4^n$ (\S\ref{sec:belnap}).
  \item \texttt{Imscribing.Millennium.ZaunerEmbeddingEquivalence}: the main equivalence and its corollaries
    (\S\ref{sec:equiv}).
  \item \texttt{Imscribing.Millennium.SIC\_D12\_ExistenceRing} and \texttt{SIC\_D12\_Embedding}:
    the exact $d = 12$ fiducial ring, the \texttt{native\_decide} identity planks, and the embedding
    that proves \texttt{SICPOVM\_Exists 12} unconditionally (\S\ref{sec:d12}).
  \item \texttt{Imscribing.Millennium.ZaunerTransportMap}: the complete, sorry-free $d = 2048$ transport:
    the forward map $\varphi$, the reduction $\psi$ with $\psi\circ\varphi=\mathrm{id}$, the conditional
    existence via Stark, and the character obstruction with its dialetheic $B$-state (\S\ref{sec:d2048}).
\end{itemize}

The structural chain is:

\[
\begin{array}{c}
\underbracket{\text{Belnap multilattice}}_{\text{sorry-free}} \\[3pt]
\big\downarrow\;{\scriptstyle \texttt{rfl}} \\[3pt]
[\,\texttt{HilbertEmbeddingExists} \;\leftrightarrow\; \texttt{ZaunerConjectureAtPow2}\,] \\[3pt]
\big\downarrow\;{\scriptstyle \text{axiom: \texttt{stark\_equivalence}}} \\[3pt]
\texttt{SICPOVM\_Exists}\,d \\[3pt]
\big\downarrow\;{\scriptstyle \text{axioms: \texttt{equiangular} + \texttt{norm}}} \\[3pt]
\texttt{MixedSignatureStarkConjecture} \\[3pt]
\big\downarrow\;{\scriptstyle \text{axiom: \texttt{zauner\_correspondence}}} \\[3pt]
\text{Stark units in }\Kd \;=\; \text{Hilbert\,12 for }\mathbb{Q}(\sqrt{(d-3)(d+1)})
\end{array}
\]

Each arrow is either proved or honestly axiomatized; the axiomatized ones carry the exact mathematical
content of what remains open.  
For $d = 12$ the terminal existence node is already stronger than the chain that predicted it:
\texttt{SICPOVM\_Exists 12} is now a proved theorem, with no axiom standing in front of it (\S\ref{sec:d12}).

The two closed dimensions sit at opposite ends of the same construction.  
At $d = 12$ the existence node itself is discharged.  
At $d = 2048$ the existence node remains the open Zauner conjecture.  
But everything around it is formalized and sorry-free: the transport $\varphi$, its reduction $\psi$,
the conditional reduction to Stark, and the character obstruction that forces the representation to be
genuinely complex.  
We carry the open node and its obstruction together as a Belnap $B$-state, refusing to resolve them
prematurely to one polarity.  
That is the honest structural status of an open problem whose skeleton is already proved.

%──────────────────────────────────────────────────────────────────────────────
\section{Weyl--Heisenberg Operators and the SIC-POVM Definition}
\label{sec:wh}
%──────────────────────────────────────────────────────────────────────────────

\subsection{The Weyl--Heisenberg group in dimension $d$}

The displacement operators in $\mathbb{C}^d$ are generated by a shift and a phase.

\begin{definition}[Shift and phase operators]
Let $\omega_d = e^{2\pi i/d}$. Define
\begin{align*}
  (X_d\,v)(k) &= v(k-1 \bmod d) \qquad (\text{shift}),\\
  (Z_d\,v)(k) &= \omega_d^k \cdot v(k) \qquad (\text{phase}).
\end{align*}
The general displacement is $D_{a,b,t} = \omega_d^t X_d^a Z_d^b$.
\end{definition}

In \Lean{} (\texttt{SIC\_POVM\_Stark.lean}, lines 26--43) they are realized concretely:

\begin{lstlisting}
def omega_d (d : N) : C := exp (2 * Real.pi * Complex.I / d)
def X_d (d : N) (v : Fin d -> C) (k : Fin d) : C :=
  v <<(k.val - 1) % d, Nat.mod_lt _ k.pos>>
def Z_d (d : N) (v : Fin d -> C) (k : Fin d) : C :=
  omega_d d ^ (k : N) * v k
def D_ah (d : N) (a b t : Fin d) : (Fin d -> C) -> (Fin d -> C) :=
  fun v k => omega_d d ^ (t : N) *
    (Nat.iterate (X_d d) (a : N) (Nat.iterate (Z_d d) (b : N) v)) k
\end{lstlisting}

These are genuine operators, not placeholder types.

\subsection{The SIC-POVM condition}

\begin{definition}[IsSICPOVM]
A vector $\psi\in\mathbb{C}^d$ is a \emph{SIC-POVM fiducial} if
\begin{enumerate}[label=(\roman*)]
  \item $\langle\psi,\psi\rangle = 1$ \quad (unit norm), and
  \item $(d+1)\,|\langle\psi, D_{a,b,0}\,\psi\rangle|^2 = 1$ for all
    $(a,b)\neq(0,0)$ \quad (equiangularity).
\end{enumerate}
\end{definition}

\begin{lstlisting}
structure IsSICPOVM (d : N) [NeZero d] (fiducial : Fin d -> C) : Prop where
  norm_eq    : wh_normSq d fiducial = 1
  equiangular : forall (a b : Fin d), (a, b) != (0, 0) ->
    ((d : R) + 1) * ||wh_inner d fiducial (D_ah d a b 0 fiducial)||^2 = 1
\end{lstlisting}

The Frobenius inner product $\langle v, w\rangle = \sum_k v_k \overline{w_k}$ is implemented
as \texttt{wh\_inner}.  
This is the standard unit-norm convention.  
An earlier draft paired $\langle\psi,\psi\rangle = d$ with $|\langle\psi,D\psi\rangle| = 1$;
those two conditions cannot hold simultaneously because the WH displacements form an orthogonal
operator basis, forcing $\sum_{a,b}|\langle\psi,D_{a,b}\psi\rangle|^2 = d\,\|\psi\|^4$, which
no vector can reconcile with the old pair.  
The unit statement above is the sound one: $\|\psi\|^2 = 1$ and every non-identity overlap has
$|\langle\psi,D_{a,b}\psi\rangle|^2 = 1/(d+1)$, written rationally to keep the condition decidable.  
Both the $d=12$ ring identities and the $d=2048$ transport machine-check this convention, and the
sanity relation $1 + (d^2-1)/(d+1) = d = d\,\|\psi\|^4$ holds.

%──────────────────────────────────────────────────────────────────────────────
\section{Arithmetic Geometry: the Stark Conjecture}
\label{sec:stark}
%──────────────────────────────────────────────────────────────────────────────

\subsection{The discriminant and base field}

\begin{definition}[Base field]
For $d \geq 2$, let $m_d = (d-3)(d+1)$ and $F_d = \mathbb{Q}(\sqrt{m_d})$.
\end{definition}

The \Lean{} definition \texttt{m\_d (d : N) : Z := ((d : Z) - 3) * ((d : Z) + 1)} is exactly the
standard SIC discriminant of Appleby, Flammia, McConnell, Yard and collaborators
\cite{Appleby2013,Appleby2017,AFMY2020}; its squarefree part fixes the field discriminant of $F_d$.
For $d \geq 4$ it is positive, so $F_d$ is real quadratic; the dimensions $d \in \{2,3\}$ are degenerate
($m_2 = -3$, $m_3 = 0$) and carry no arithmetic obstruction.  
The dimension $d$ \emph{selects} the arithmetic field: the SIC-POVM problem in dimension $d$ is the
Stark problem over $F_d$.  
The conjecture is therefore not a single universal statement but a family of arithmetic problems,
one per real quadratic field, each independently open or closed.

\begin{remark}
For $d=4$: $m_4 = 5$, $F_4 = \mathbb{Q}(\sqrt{5})$, the golden-ratio field of the Hesse SIC.  
For $d=5$: $m_5 = 12$, $F_5 = \mathbb{Q}(\sqrt{3})$; for $d=7$: $m_7 = 32$, $F_7 = \mathbb{Q}(\sqrt{2})$.  
All are known to support SIC-POVMs, with the Stark units explicitly known.  
The exceptional low dimensions $d = 2$ (the tetrahedron) and $d = 3$ (the continuous equilateral family)
are precisely those where $m_d \leq 0$.
\end{remark}

\subsection{The mixed-signature Stark conjecture}

The mixed-signature Stark conjecture asserts the existence of a unit $\varepsilon_d \in K_d^\times$
with controlled embedding absolute values, satisfying the regulator condition
$\log|\varepsilon_d^\sigma| = L'(0,\chi^\sigma)$ for each $L$-function twist.

\begin{axiom_env}[MixedSignatureStarkConjecture]
\label{ax:stark}
For $d \geq 4$ with $m_d$ non-square:
\[
\forall\, i,j \in \{0,\ldots,d-1\}:\quad
\|\sigma_i(\varepsilon_d)\| \leq \tfrac{1}{d} + 1
\quad\text{and}\quad
\forall\, \tau \in \mathrm{Gal}(K_d/F_d):\; \sigma_j(\tau\cdot\varepsilon_d) \neq 0
\]
\end{axiom_env}

This is an open problem in analytic number theory.  
In \Lean{} it is marked as \texttt{def MixedSignatureStarkConjecture}, not \texttt{axiom};
it is a defined proposition whose truth value is unknown.

\subsection{Constructing the fiducial from the Stark unit}

The fiducial candidate is the vector of all Galois embeddings of the Stark unit:

\[
v_d(k) = \sigma_k(\varepsilon_d),\quad k = 0,\ldots,d-1
\]

\begin{lstlisting}
def fiducial_from_stark (d : N) (hd : 4 <= d) (hns : not IsSquare (m_d d)) :
    Fin d -> C :=
  fun k => (Embeddings d hd hns k) (StarkUnit d hd hns)

def normalize_fiducial ... :=
  fun k => (Real.sqrt (d : R))^(-1) * fiducial_from_stark d hd hns k
\end{lstlisting}

\subsection{The Galois--Zauner correspondence}

\begin{axiom_env}[zauner\_correspondence]
\label{ax:zauner-corr}
The absolute value of the WH orbit inner product is controlled by the order-3 Zauner automorphism
$\zeta_d \in \mathrm{Gal}(K_d/F_d)$ acting on the Stark unit:
\[
\big|\langle v_d,\, D_{a,b}\,v_d\rangle\big|
= \big|\langle v_d, v_d\rangle\big| \cdot
\big|\overline{\sigma_0(\zeta_d \cdot \varepsilon_d)}\big|
\]
\end{axiom_env}

This axiom captures the Appleby--Yadsan-Appleby--Zauner result \cite{Appleby2013}: the SIC overlap
is an algebraic function of the Stark unit under the Zauner automorphism.
It is formalized as a named \texttt{axiom} with the exact mathematical content stated in the type.

\subsection{Honest gap markers}
\label{sec:gaps}

Two further axioms carry the open arithmetic geometry:

\begin{axiom_env}[equiangular\_from\_stark\_axiom, norm\_of\_normalized\_axiom]
\label{ax:gap}
Given the mixed-signature Stark conjecture~\ref{ax:stark}, the normalized fiducial satisfies:
\begin{itemize}
  \item $(d+1)\,|\langle \tilde{v}_d, D_{a,b}\,\tilde{v}_d\rangle|^2 = 1$ for all $(a,b)\neq(0,0)$,
  \item $\langle\tilde{v}_d, \tilde{v}_d\rangle = 1$ \quad (unit norm).
\end{itemize}
These are exactly the two clauses of \texttt{IsSICPOVM} in the unit convention of \S\ref{sec:wh},
as the \Lean{} statements \texttt{equiangular\_from\_stark\_axiom} and
\texttt{norm\_of\_normalized\_axiom} read verbatim.
\emph{Gap:} the full arithmetic geometry of Stark units acting on WH frames: the translation from the
Stark embedding bounds to the exact equiangularity and norm conditions is the open mathematical content.
\end{axiom_env}

Both axioms carry the same source comment:
\begin{quotation}
\noindent\textit{``the gap between the Stark conjecture and these norm statements is the open content;
it requires the full arithmetic geometry of Stark units acting on WH frames.''}
\end{quotation}

This is not evasion. It is the precise location of what remains to be done.

\subsection{The conditional existence theorem}

\begin{theorem}[sic\_povm\_exists\_via\_stark]
\label{thm:stark-sic}
Assume the mixed-signature Stark conjecture and the two gap axioms (\ref{ax:gap}).
Then $\texttt{SICPOVM\_Exists}(d)$ holds for all $d \geq 4$ with $m_d$ non-square.
\end{theorem}

\begin{lstlisting}
theorem sic_povm_exists_via_stark
    (d : N) [NeZero d] (hd : 4 <= d) (hns : not IsSquare (m_d d))
    (sc : MixedSignatureStarkConjecture d hd hns) :
    SICPOVM_Exists d := by
  use normalize_fiducial d hd hns
  exact { norm_eq := norm_of_normalized d hd hns sc,
          equiangular := equiangular_from_stark d hd hns sc }
\end{lstlisting}

%──────────────────────────────────────────────────────────────────────────────
\section{The Belnap Multilattice}
\label{sec:belnap}
%──────────────────────────────────────────────────────────────────────────────

\subsection{The Belnap four-valued lattice $\BFour$}

The Belnap--Dunn four-valued logic \cite{Belnap1977} has truth values $\{N, T, F, B\}$
(neither, true, false, both) ordered by the bilattice structure.  
The key property needed here is the dialetheic fixed point:
\[
\neg B = B \qquad (\texttt{bnot B = B}).
\]
$B$ absorbs negation.

\subsection{The $n$-qubit Belnap state and WH action}

An $n$-qubit Belnap state is a pair $(\text{val}, \text{phase}) \in \BFour \times \mathbb{Z}_2$.
The all-$B$ all-zero state $B^{\otimes n}$ is the fiducial.  
The WH displacement acts componentwise: an amplitude displacement $a_i = 1$ applies $\neg$ to
position $i$; because $\neg B = B$, amplitude displacements are invisible on $B^{\otimes n}$.

\begin{theorem}[whOrbit\_card\_eq\_pow2]
\label{thm:orbit-2n}
The product-lattice WH orbit of $B^{\otimes n}$ has cardinality $2^n$:
\[
|\mathcal{O}_{B^{\otimes n}}| = 2^n
\]
\end{theorem}

\begin{proof}
Amplitude displacements fix $B^{\otimes n}$ (since $\neg B = B$).  
Displaced states differ only in the phase word $p \in (\mathbb{Z}_2)^n$.  
The phase embedding $p \mapsto (B, p)^{\otimes n}$ is injective and its image equals the orbit, so
$|\mathcal{O}| = |(\mathbb{Z}_2)^n| = 2^n$.  
\textbf{Proved in \Lean{}, 0 sorries.}
\end{proof}

\begin{corollary}[product\_lattice\_orbit\_is\_insufficient]
\label{cor:gap}
For $n \geq 1$: $|\mathcal{O}_{B^{\otimes n}}| = 2^n < 4^n = d^2$.  
The product-lattice orbit is insufficient for a SIC-POVM.
\end{corollary}

This is the \emph{orbit gap}: the structural gap that the multilattice axioms close.

\subsection{The multilattice axioms}

An extended type $\texttt{BelnapML}(n)$ is postulated to support amplitude-distinguishable
displacements. Four axioms specify its properties.

\begin{axiom_env}[Ax-PROJ]
The multilattice fiducial projects to $B^{\otimes n}$ in the product lattice.
\end{axiom_env}

\begin{axiom_env}[Ax-FREE]
The WH action on the multilattice fiducial produces $4^n$ distinct states:
$g \neq h \implies g \cdot B^{\otimes n} \neq h \cdot B^{\otimes n}$.
\end{axiom_env}

\begin{axiom_env}[Ax-EQUI]
WH-displaced fiducials have constant pairwise overlap:
$\exists\, k,\; \forall\, g \neq h:\; \langle g \cdot B^{\otimes n}, h \cdot B^{\otimes n}\rangle = k$.
This is the SIC-POVM equiangularity condition in the Belnap metric.
\end{axiom_env}

\begin{axiom_env}[Ax-COST]
The WH SIC measurement yields $\texttt{belnapCost} = 2 \times \texttt{period}$ for any coprime pair
$(N, a)$.
\end{axiom_env}

\begin{proposition}[belnap\_structural\_skeleton]
\label{prop:skeleton}
For every $n \geq 1$, the Belnap multilattice satisfies: the orbit size is $4^n$ (by Ax-FREE);
the meet-identity $\bigwedge_{x}\, \mathrm{wordMeet} = x$ holds; distinct group elements produce
distinct states; and the join-equiangularity condition
$\mathrm{frobInner}(B^{\otimes n}, g \cdot B^{\otimes n}) = 2n$ holds for all $g$.
All of these follow from \texttt{sic\_povm\_belnap\_unconditional} (0 sorries).
\end{proposition}

\subsection{The group-theoretic bifurcation}

The equivalence between the Belnap multilattice and the Zauner conjecture is nontrivial because the
index groups differ:
\[
\WH(2)^n \cong (\mathbb{Z}_2)^{2n}
\qquad\text{vs}\qquad
\WH(2^n) \cong \mathbb{Z}_{2^n} \times \mathbb{Z}_{2^n}.
\]
The characters of the elementary abelian group $\WH(2)^n$ are $\pm 1$-valued.
The module's commentary records a parity-gated conjecture about whether the product structure can
nevertheless meet the SIC overlap magnitude $1/\sqrt{2^n+1}$: affirmed at $n \in \{1,3\}$,
conjectured for odd $n$, and conjectured to fail for even $n$; only the cardinality statement
below is proved.
In either parity, realizing the metric SIC requires the lift $\WH(2)^n \to \WH(2^n)$, and this
lift \emph{is} the Zauner conjecture.

\begin{theorem}[group\_bifurcation\_lemma]
$|\texttt{WHIdx}(n)| = 4^n$ (the index group has the right cardinality).
\end{theorem}

%──────────────────────────────────────────────────────────────────────────────
\section{The Main Equivalence}
\label{sec:equiv}
%──────────────────────────────────────────────────────────────────────────────

\subsection{Definitions}

\begin{definition}
\[
\texttt{ZaunerConjectureAtPow2}(n) \;:=\; \texttt{SICPOVM\_Exists}(2^n)
\]
\[
\texttt{HilbertEmbeddingExists}(n) \;:=\; \texttt{SICPOVM\_Exists}(2^n)
\]
\end{definition}

Both are defined as the same proposition: the existence of a WH$(2^n)$-covariant SIC-POVM fiducial
in $\mathbb{C}^{2^n}$.  
The Belnap multilattice provides the structural content (orbit size $4^n$, equiangularity, Frobenius closure);
\texttt{SICPOVM\_Exists} provides the metric realizability condition.

\subsection{The equivalence is proved by \texttt{rfl}}

\begin{theorem}[hilbert\_embedding\_equiv\_zauner]
\label{thm:main-equiv}
For every $n \geq 1$:
\[
\texttt{HilbertEmbeddingExists}(n) \;\longleftrightarrow\; \texttt{ZaunerConjectureAtPow2}(n)
\]
\end{theorem}

\begin{lstlisting}
theorem hilbert_embedding_equiv_zauner (n : N) [NeZero (2 ^ n)] :
    HilbertEmbeddingExists n <-> ZaunerConjectureAtPow2 n := by
  rfl
\end{lstlisting}

The proof is one token.  
This is not a triviality; it is a structural identification.  
Both definitions reduce to \texttt{SICPOVM\_Exists (2\^{}n)}: the Belnap multilattice embedding
into $\mathbb{C}^{2^n}$ and the Zauner conjecture at that dimension are, definitionally, the same
open problem.

The force of the theorem lies in what it identifies, not in the proof term.  
The Belnap multilattice (22 theorems, 0 sorries) proves \emph{all} SIC structural axioms in the
Belnap metric unconditionally.  
The \texttt{rfl} says: the remaining gap (can this structure be represented in $\mathbb{C}^{2^n}$
with the standard WH$(2^n)$ action?) is exactly what the Zauner conjecture asks.

\subsection{Corollary: closing one level closes all three}

\begin{corollary}
\label{cor:chain}
If the mixed-signature Stark conjecture holds for $\Kd = \mathbb{Q}(\sqrt{(d-3)(d+1)})$,
then:
\begin{enumerate}[label=(\roman*)]
  \item \texttt{sic\_povm\_exists\_via\_stark} closes:
    $\texttt{SICPOVM\_Exists}(d)$ holds.
  \item \texttt{hilbert\_embedding\_equiv\_zauner} closes: the Belnap multilattice
    embeds into $\mathbb{C}^d$.
  \item The ray class field $\Kd$ is generated by the fiducial coordinates
    (Hilbert's Twelfth Problem for $F_d$).
\end{enumerate}
\end{corollary}

The paraconsistent quantum information structure and the arithmetic geometry of real quadratic fields
are formally identified as \emph{the same problem}.

%──────────────────────────────────────────────────────────────────────────────
\section{Hilbert's Twelfth Problem}
\label{sec:hilbert12}
%──────────────────────────────────────────────────────────────────────────────

Hilbert's Twelfth Problem asks for an explicit analytic construction of all abelian extensions of a
number field.  
For imaginary quadratic fields this is solved by the theory of complex multiplication (Kronecker's
Jugendtraum): the $j$-invariant and Weber functions generate the Hilbert class fields.

The real quadratic case is entirely open.  
Stark's approach \cite{Stark1980} via $L$-functions at $s=0$ succeeds for totally imaginary abelian
extensions but runs into the \emph{mixed signature} problem for real quadratic base fields:
the leading coefficient of the Taylor expansion of $L(s,\chi)$ at $s=0$ involves the regulator of
the Stark unit, and controlling its complex embeddings requires additional non-vanishing conditions
that have not been established in general.

The formalization connects these levels through the discriminant.

\begin{remark}[Discriminant formula]
$m_d = (d-3)(d+1)$ selects the base field $F_d = \mathbb{Q}(\sqrt{m_d})$; the field discriminant
is that of the squarefree part of $m_d$ (for $d = 12$, $m_{12} = 9\cdot 13$ and
$\operatorname{disc} F_{12} = 13$).
For $d \geq 4$: $m_d > 0$ (real quadratic).
The integer $m_d$ is not a perfect square for generic $d$, and the non-square condition
\texttt{hns} appears throughout the formalization as a hypothesis, not a theorem.  
The explicit $d=12$ computation of \S\ref{sec:d12} realizes this concretely:
$m_{12} = 9\cdot 13$, so $F_{12} = \mathbb{Q}(\sqrt{13})$, confirmed to $1500$ significant digits.
\end{remark}

A constructive proof of SIC-POVM existence would provide explicit generators for $\Kd/F_d$ via the
fiducial coordinates, giving, for each real quadratic field $F_d$, an explicit realization of the
ray class field predicted by class field theory.  
That would constitute a proof of Hilbert's Twelfth Problem in the real quadratic case.

%──────────────────────────────────────────────────────────────────────────────
\section{The \texorpdfstring{$d=12$}{d=12} Existence Theorem}
\label{sec:d12}
%──────────────────────────────────────────────────────────────────────────────

The conditional construction of \S\ref{sec:stark} predicts that the fiducial coordinates generate
a specific ray class field.  
For $d=12$ this chapter of the paper closes completely.  
We exhibit the predicted field as an explicit, computable object, verify its arithmetic by machine,
construct an exact fiducial whose twelve coordinates live in an explicitly presented algebra over
$\mathbb{Q}$, and prove in \Lean{}, with no hypothesis, that it is a SIC-POVM:
\[
\texttt{theorem crystal\_forces\_d12\_sic : SICPOVM\_Exists 12}
\]
is a sorry-free theorem of the development, resting on no axiom beyond \Lean{}'s standard foundations
and the \texttt{native\_decide} compiler trust (\S\ref{subsec:audit}).  
In earlier drafts this section was titled a \emph{witness}; the word is now too weak.  
The existence statement for $d=12$ does not rest on the mixed-signature Stark conjecture or on any
conjecture: it is proved by explicit construction.  
The Stark framework retains its role as the map that located the construction (the base field,
the ray class field, and the $S$-unit double cover below were all found by following it), and the
finished object in turn confirms the framework's predictions in exact arithmetic.  
This is a concrete instance of the Hilbert-Twelfth construction of \S\ref{sec:hilbert12}, realized
end to end.

\subsection{The base field is \texorpdfstring{$\mathbb{Q}(\sqrt{13})$}{Q(sqrt13)}}

The SIC discriminant of Appleby, Flammia, McConnell, Yard and collaborators
\cite{Appleby2013,Appleby2017,AFMY2020} is $(d-3)(d+1)$.
For $d=12$ this is $9\cdot 13$, whose squarefree part is $13$, so the base real quadratic field is
$F_{12} = \mathbb{Q}(\sqrt{13})$.  
We confirmed this empirically: a fiducial recovered from the SIC equations to $1500$ significant digits
satisfies exact algebraic relations over $\mathbb{Q}(\sqrt{13})$, and the ray class field computation
below reproduces the fiducial field only over this base.  
In particular $\sqrt{13}$, not $\sqrt{30}$, occurs throughout.

\begin{remark}[A discarded alternate]
A superficially plausible alternate discriminant $m_d = d(d-2)$ would give, for $d=12$,
$\mathbb{Q}(\sqrt{120}) = \mathbb{Q}(\sqrt{30})$.  
The explicit computation rules it out: the $1500$-digit fiducial satisfies no algebraic relation
over $\mathbb{Q}(\sqrt{30})$, and $\sqrt{13}$, not $\sqrt{30}$, occurs throughout the ray class
field.  
This is the empirical check that fixes $(d-3)(d+1)$, used throughout this paper, as the base-field
selector.
\end{remark}

\subsection{The coordinate field}

Using PARI/GP (\texttt{bnrinit}, \texttt{bnrclassfield}) we find that the $d=12$ fiducial coordinates
generate the ray class field $K_{12}$ of $F_{12} = \mathbb{Q}(\sqrt{13})$ of conductor
$\mathfrak{f} = (3d)\,\infty_1\infty_2 = (36)\,\infty_1\infty_2$, with both infinite places ramified.  
Its ray class group is
\[
\mathrm{Cl}_{\mathfrak{f}}(F_{12}) \;\cong\;
\mathbb{Z}/6 \times \mathbb{Z}/6 \times \mathbb{Z}/2 \times \mathbb{Z}/2,
\qquad |\mathrm{Cl}_{\mathfrak{f}}| = 144,
\]
so $[K_{12} : \mathbb{Q}] = 2\cdot 144 = 288$.  
The moduli $|v_{12}(k)|^2$ are \emph{real}, and they generate the index-$2$ real subfield, of
degree $16$, of the conductor-$(12)\,\infty_1\infty_2$ ray class field of degree $32$, which
carries $i$.
That subfield is fixed by complex conjugation at the fiducial embedding but is \emph{not}
totally real: its signature is $(8,4)$.
The individual moduli split by index: $N_0, N_3, N_6, N_9$ have minimal polynomials of degree
$4$ over $\mathbb{Q}$ and the remaining eight have degree $8$.
Two independent primitive elements, $\sum_k (k{+}1)\,|v_{12}(k)|^2$ and
$\sum_k (2k{+}1)\,|v_{12}(k)|^2$, each have minimal polynomial of degree exactly $16$, fixing the
moduli field.
As \S\ref{subsec:doublecover} shows, the full coordinates generate a proper extension of $K_{12}$.

\subsection{Explicit tower decomposition}

The class field decomposes (via \texttt{bnrclassfield}) as a compositum of six cyclic pieces over
$F_{12}$, each with a small defining polynomial; the six pieces are collected in Table~\ref{tab:tower}.
\begin{table}[h]
\centering
\begin{tabular}{@{}lll@{}}
\toprule
\textbf{Piece} & \textbf{Defining polynomial over $\mathbb{Q}(\sqrt{13})$} & \textbf{Generator} \\
\midrule
$p_1$ & $x^2 + 1$ & $i$ \\
$p_2$ & $x^2 - (\sqrt{13}+5)/2$ & $\sqrt{(5+\sqrt{13})/2}$ \\
$p_3$ & $x^2 - (\sqrt{13}-1)/2$ & $\sqrt{(\sqrt{13}-1)/2}$ \\
$p_4$ & $x^2 + (\sqrt{13}+1)/2$ & $\sqrt{-(\sqrt{13}+1)/2}$ \\
$p_5$ & $x^3 - 3x - 1$ & real cubic of $\mathbb{Q}(\zeta_9)$ \\
$p_6$ & $x^3 + (33\sqrt{13}-123)x + (215 - 59\sqrt{13})$ & SIC cubic \\
\bottomrule
\end{tabular}
\caption{Explicit tower decomposition of the class field into six cyclic pieces over $\mathbb{Q}(\sqrt{13})$.}
\label{tab:tower}
\end{table}
Four quadratic layers and two cyclic cubics give $2^5\cdot 3^2 = 288$.  
The conductor $36 = 4\cdot 9$ factors transparently: $i$ contributes the $4$, the $\zeta_9$ cubic
contributes the $9$, over the $\sqrt{13}$ arithmetic.

\subsection{The coordinate field is a ramified \texorpdfstring{$S$}{S}-unit double cover}
\label{subsec:doublecover}

The construction $v_{12}(k) = \sigma_k(\varepsilon_{12})$ of \S\ref{sec:stark} would place every
coordinate in $K_{12}$, hence at degree dividing $288$.  
A high-precision check refutes this for the crux coordinates and, in doing so, sharpens the construction.
We recovered the fiducial to $8000$ significant digits and tested each coordinate by a
\emph{deep-vanish floor} criterion: a genuine minimal polynomial vanishes to full precision and then
plateaus in degree, while a low-precision artifact creeps without flooring.  
We find $v_{12}(1)$ irreducible of degree $64 = 2^6$.  
Since $2^6 \nmid 288 = 2^5\cdot 3^2$, \emph{$v_{12}(1)$ does not lie in $K_{12}$}; the twelve
coordinates generate a proper extension.

The structure is a ramified quadratic double cover.  
Write $v_{12}(k) = |v_{12}(k)|\,u_k$ with unit phase $u_k = v_{12}(k)/|v_{12}(k)|$.  
The phase $u_1$ has degree $32 \mid 288$ and lies in $K_{12}$, and the square
$v_{12}(1)^2 = N_1\,u_1^2$ (where $N_1 = |v_{12}(1)|^2$) is irreducible of degree $32$ and also lies
in $K_{12}$.  
Hence $v_{12}(1) = \sqrt{v_{12}(1)^2}$ is a quadratic extension of a $K_{12}$ element, and the degree
$2^6$ is precisely one factor of $2$ above the $2^5$ part of $K_{12}$: the square root of the modulus.

That square root is a unit at the ramified primes.  
Modulo squares the ramifying class is $[N_1]$, the phase square $u_1^2$ dropping out.  
The minimal polynomial of $N_1$ over $\mathbb{Q}$ has degree $8$, with field norm
\[
\mathrm{N}_{\mathbb{Q}}(N_1) \;=\; \frac{9}{9475854336} \;=\; \frac{1}{32448^{\,2}},
\qquad 32448 = 2^6\cdot 3\cdot 13^2,
\]
a perfect square supported only on the primes $\{2,3,13\}$ ramified in $K_{12}$ (field discriminant
$2^{12}\,3^2\,13^4$).  
Thus $K_{12}(\sqrt{N_1})/K_{12}$ ramifies only where $K_{12}$ already ramifies: $[N_1]$ is an
$S$-unit class with $S = \{2,3,13\}$, and
\[
v_{12}(1) \;=\; \sqrt{u}\,\cdot\,(\text{element of } K_{12}),
\qquad u \text{ an } S\text{-unit.}
\]
The exact-embedding construction of \S\ref{sec:stark} is therefore refined to
$v_{12}(k) = \sigma_k(\sqrt{\varepsilon})$ for an $S$-unit $\varepsilon$: the coordinate is the
embedding of the \emph{square root} of the Stark datum, not the datum itself, which resolves the
degree count $64 = 2\cdot 32$.  
This is the arithmetic shadow of the Frobenius self-pairing $\frob$ of \S\ref{sec:belnap}: the modulus
$N_k = v_{12}(k)\,\overline{v_{12}(k)}$ is the idempotent collapse of a coordinate against its
mirror into the real subfield, and the coordinate is that collapse's square-root double cover.  
The base $F_{12} = \mathbb{Q}(\sqrt{13})$ has fundamental unit $\varepsilon_{13} = (3+\sqrt{13})/2$
(regulator $1.19476\ldots$); the exact representative is pinned in \S\ref{subsec:ring}, where the
whole coordinate system is presented as one finite ring.

\subsection{The existence ring}
\label{subsec:ring}

The double cover says where the coordinates live; the finished construction presents that home as a
single finite object.  
Let $K_{16} = \mathbb{Q}[g]/(p_r)$ be the degree-$16$ moduli field (real at the fiducial
embedding, signature $(8,4)$), with
$p_r = x^{16}-10x^{14}+40x^{12}-90x^{10}+126x^{8}-96x^{6}+25x^{4}+2x^{2}+1$, in which all twelve
squared moduli $N_k = |v_{12}(k)|^2$ are exact rational vectors.  
Over $K_{16}$ we adjoin four magnitude square roots $s_0, s_1, s_3, s_9$ (one per independent cover
class of \S\ref{subsec:doublecover}), the imaginary unit $i$, one real quadratic layer $c_5$, and
one phase generator $u_1$.  
The result is a commutative ring with involution
\[
R \;=\; K_{16}(s_0, s_1, s_3, s_9,\; i,\; c_5,\; u_1),
\qquad \dim_{\mathbb{Q}} R \;=\; 2048 .
\]
Two collapse relations, found by an exhaustive branch audit of the formal engine against the
$1500$-digit witness, cut the tower to this size.  
First, the phase $u_1$ is \emph{quadratic} over $K_{16}(i)$: its square is the half-angle datum
$(c_2 + i\,s_2)/2$ with $c_2, s_2 \in K_{16}$, so the expected octic phase layer is four times
smaller than its naive degree count.  
Second, the fifth cover collapses into the $c_5$ layer through the exact relation
$N_1 N_5 (A_5^2 - 4B_5) = \rho^2$ with $\rho \in K_{16}$, which also places $\sqrt{3}$ and
$\zeta_{12}$ inside $K_{16}(c_5, i)$.

All twelve coordinates are named elements of $R$, and the SIC conditions hold in $R$ \emph{exactly}:
the norm identity $\sum_k N_k = 1$, the coordinate moduli $\bar z_k z_k = N_k$, and all $143$ overlap
identities $O_{a,b}\,\overline{O_{a,b}} = \tfrac{1}{13}$ are theorems of
\texttt{SIC\_D12\_ExistenceRing.lean}.  
Each is closed by \texttt{native\_decide} over exact rational arithmetic (a sparse associative-list
engine on $7$-bit monomial keys; no floating point appears anywhere in the file).  
The Lean source is generated by a script that re-executes an independent exact-fraction engine as a
gate and emits its vectors verbatim, so the formal data has a single source of truth.  
Earlier computable models (a quadratic-extension stack for the moduli subfield and a flat degree-$288$
reduction engine in which $\theta^{288}$ reduces under \texttt{native\_decide}) remain in the
development as the route that located this presentation.

The structural payoff of working in $R$: the identities are ring identities, so
\emph{every} star-compatible ring homomorphism $R \to \mathbb{C}$ sends the twelve named elements
to a SIC fiducial.  
Existence in $\mathbb{C}^{12}$ reduces to exhibiting one homomorphism.

\subsection{The embedding and the audit}
\label{subsec:audit}

\texttt{SIC\_D12\_Embedding.lean} builds the homomorphism $\varphi : R \to \mathbb{C}$ layer by layer
and proves it is a star-ring map on the canonical monomial domain.  
The base point $g \mapsto g_0$ is the real root of $p_r$ in the rational bracket $(-2.009,\,-2.008)$,
produced by the intermediate value theorem from exact sign changes at rational endpoints (verified by
\texttt{native\_decide} in $\mathbb{Q}$, then lifted to $\mathbb{R}$).  
The four cover generators map to real square roots of the evaluated moduli; their nonnegativity, the
one genuinely analytic obligation, is discharged by exact divided-difference certificates: for each
polynomial value, the quotient $q(x) = (p(x)-p(m))/(x-m)$ is computed exactly in $\mathbb{Q}[x]$,
giving the bound $p(x) \ge p(m) - \tfrac{w}{2}\sum_i |q_i|\,B^i$ on the bracket by pure rational
arithmetic.  
No interval arithmetic and no floating point enter the proof.  
The layer $c_5$ maps to the real quadratic root selected by a positive-discriminant certificate, and
$u_1$ is reconstructed by half-angle, $u_1 = \sqrt{(1+x)/2} + i\,\sqrt{(1-x)/2}$, which makes
conjugation compatibility a theorem rather than a branch convention.

The transfer is then mechanical in the best sense.  
The norm half of the SIC condition is the frozen ring identity $\sum_k N_k = 1$ pushed through
$\varphi$.  
For equiangularity, the entire ring side ($|\varphi(O_{a,b})|^2 = \tfrac{1}{13}$ for all $143$ pairs)
follows from the \texttt{native\_decide} identities through the homomorphism; one bridge lemma matches
the Weyl--Heisenberg displacement iterates of \S\ref{sec:wh} against the ring overlaps.  
The bridge pins $Z = \varphi(\zeta)$ by exact computation of its imaginary part
($\operatorname{Im} Z = \tfrac12$) and its twelfth power ($Z^{12} = 1$), so $Z \in \{\omega, \omega^5\}$;
both branches are handled by the reindexing $b' = (m\,b) \bmod 12$, a bijection on the $143$ nontrivial
pairs because $m \in \{1,5\}$ is a unit modulo $12$.

\begin{lstlisting}
theorem d12_sic_exists : IsSICPOVM 12 psi

theorem crystal_forces_d12_sic : SICPOVM_Exists 12
\end{lstlisting}

The axiom audit closes the account.  
\texttt{\#print axioms} on \texttt{crystal\_forces\_d12\_sic} reports exactly \texttt{propext},
\texttt{Classical.choice}, \texttt{Quot.sound}, \texttt{Lean.ofReduceBool}, and
\texttt{Lean.trustCompiler}: \Lean{}'s standard three plus the compiler trust that
\texttt{native\_decide} always carries.  
Not a single axiom of this development appears in the list; not the Stark axioms of \S\ref{sec:stark},
not any shadow assumption.  
The declaration \texttt{crystal\_forces\_d12\_sic} entered the development as an \texttt{axiom};
it is now a theorem, and the axiom is deleted.

\subsection{Scope}

What is established unconditionally: a WH-covariant SIC-POVM exists in dimension $12$, by explicit
construction, machine-checked from the ring presentation down to the kernel.  
Exact $d=12$ fiducials were known in the literature; the formal, audited existence proof is, to our
knowledge, the first in any proof assistant, in any dimension.  
What is established about the arithmetic: the construction confirms the Stark-theoretic predictions
in exact form (base field $\mathbb{Q}(\sqrt{13})$, ray class field of conductor $(36)\,\infty_1\infty_2$,
coordinates on the ramified $S$-unit double cover), so the $d=12$ object is a fully verified instance
of the Hilbert-Twelfth construction of \S\ref{sec:hilbert12}.  
What is \emph{not} claimed: no general-$d$ statement follows from the construction; the mixed-signature
Stark conjecture is not proved, and the general chain of \S\ref{sec:stark} remains exactly as
axiomatized in \S\ref{sec:gaps}.

%──────────────────────────────────────────────────────────────────────────────
\section{The \texorpdfstring{$d=2048$}{d=2048} Zauner Transport}
\label{sec:d2048}
%──────────────────────────────────────────────────────────────────────────────

Dimension $d = 2048 = 2^{11}$ is the frontier case of the $d = 2^n$ family.  
Unlike $d = 12$, its unconditional existence is not settled here: it remains the open Zauner conjecture,
identified with the mixed-signature Stark conjecture for $F_{2048}$ by Theorem~\ref{thm:main-equiv}.  
What is settled, sorry-free in \texttt{ZaunerTransportMap.lean}, is the entire transport apparatus that
surrounds that one open node.  
We prove the unconditional Belnap skeleton, the forward transport and its reduction, the conditional
reduction to Stark, and a genuine obstruction to real representations; and we carry the open existence
node together with that obstruction as a Belnap $B$-state rather than resolving it prematurely.

\subsection{The problem space and the unconditional skeleton}

Write $\texttt{zauner\_d} = 2048$ and $\texttt{zauner\_n} = 11$, with the fiducial $B^{\oplus 11}$
the all-$B$ Belnap word of length $11$.  
The structural skeleton is inherited from \texttt{sic\_povm\_belnap\_unconditional} of \S\ref{sec:belnap},
instantiated at $n = 11$ and proved with zero \texttt{sorry}s:
\[
\begin{gathered}
|\mathcal{O}_{B^{\oplus 11}}| = 4^{11} = 4{,}194{,}304,\\
\mathrm{frobInner}\!\left(B^{\oplus 11},\, g\cdot B^{\oplus 11}\right) = 22
\quad\text{for every } g \in \WH\mathrm{Idx}(11).
\end{gathered}
\]
The orbit has the SIC cardinality $d^2$, and the join-equiangularity constant $2n = 22$ is uniform
across the orbit.  
No arithmetic hypothesis is used for either statement.

\subsection{The transport map and its reduction}

The forward transport assigns a complex amplitude to each of the $2048$ coordinates from the Hamming
weight of its bit vector.  
For $k \in \mathrm{Fin}\,2048$ with bit vector $v \in \{0,1\}^{11}$ of weight $|v|$,
\[
\varphi\!\left(B^{\oplus 11}\right)(k) = \frac{\omega_d^{\,|v|}}{\sqrt{d}},
\qquad \omega_d = e^{2\pi i/d}.
\]

\begin{theorem}[transport\_preserves\_norm\_sq]
\label{thm:transport-unit}
$\varphi(B^{\oplus 11})$ is a unit vector:
$\texttt{wh\_normSq}\,2048\,\allowbreak \varphi(B^{\oplus 11}) = 1$.
\end{theorem}

\begin{proof}
Each amplitude has modulus $1/\sqrt{d}$ because $\omega_d$ lies on the unit circle and
$\|\omega_d^{\,|v|}\| = 1$.  
Hence every coordinate contributes $|{\cdot}|^2 = 1/d$, and the $d = 2048$ coordinates sum to
$d \cdot (1/d) = 1$.  
Proved in \Lean{}.
\end{proof}

The reduction $\psi\colon \mathbb{C}^{2048} \to B^{\oplus 11}$ reads a Belnap marginal at each
register $i$ by probing two representative coordinates: index $0$ (bit $i$ clear) and index $2^i$
(bit $i$ set).  
Both carrying amplitude reads $B$; only the clear probe reads $T$; only the set probe reads $F$;
neither reads $N$.  
This is the structural content of the reverse morphism \texttt{AREV}, the read-back half of the
split/fuse pair of \S\ref{sec:belnap}.

\begin{theorem}[reduction\_recovers\_fiducial]
\label{thm:reduction}
$\psi \circ \varphi = \mathrm{id}$ on the fiducial:
$\psi\!\left(\varphi(B^{\oplus 11})\right) = B^{\oplus 11}$.
\end{theorem}

\begin{proof}
Every transport amplitude is nonzero (a unit phase times $1/\sqrt d \neq 0$), so at every register
both probes fire and the marginal reads $B$.  
The result is the all-$B$ word $B^{\oplus 11}$.  
Proved in \Lean{}.
\end{proof}

Theorems~\ref{thm:transport-unit} and~\ref{thm:reduction} exhibit $\varphi$ and $\psi$ as a
transport/reduction pair realizing the Frobenius adjoint half $\frob$ concretely in
$\mathbb{C}^{2048}$: the skeleton's structural content survives the round trip into the Hilbert
space and back.

\subsection{The existence ring and the conditional reduction}

The arithmetic geometry mirrors $d = 12$.  
The SIC discriminant is
$m_{2048} = (d-3)(d+1) = 2045\cdot 2049 = 4{,}190{,}205 = 3\cdot 5\cdot 409\cdot 683$,
squarefree, so $F_{2048} = \mathbb{Q}(\sqrt{4190205})$ is real quadratic and the non-square hypothesis
holds by \texttt{native\_decide}.  
Over it sit the ray class field $K_{2048}$, the Stark unit $\varepsilon_{2048}$, and the $2048$
complex embeddings, all instantiated from \S\ref{sec:stark} at $d = 2048$.

\begin{theorem}[zauner\_transport\_theorem]
\label{thm:d2048-cond}
If the mixed-signature Stark conjecture holds for $K_{2048}$, then
$\texttt{SICPOVM\_Exists}\,2048$.
\end{theorem}

Under the same hypothesis the structural join-equiangularity constant $22$ lifts to the metric
equiangularity $(d+1)\,|\langle\psi, D_{a,b}\psi\rangle|^2 = 1$, i.e.\
$|\langle\psi,D_{a,b}\psi\rangle|^2 = 1/2049$, reconciling the discrete skeleton with the
$\mathbb{C}^{2048}$ frame (the fuse morphism \texttt{FFUSE} of the same pair).

\subsection{The character obstruction}
\label{subsec:obstruction}

The reduction to Stark is the \emph{true} branch.  
The \emph{false} branch is a genuine obstruction to staying real.  
We prove it exactly, and state its scope precisely.

\begin{theorem}[evalf\_character\_obstruction]
\label{thm:obstruction}
For $d = 2^n$ with $n > 1$, the Weyl--Heisenberg phase character $\omega_d = e^{2\pi i/d}$ is not real:
there is no $r \in \mathbb{R}$ with $\omega_d = r$.
\end{theorem}

\begin{proof}
$\omega_d = \exp\!\big((2\pi/d)\,i\big)$ has imaginary part $\operatorname{Im}\omega_d = \sin(2\pi/d)$.  
For $d = 2^n > 2$ we have $0 < 2\pi/d < \pi$, so $\sin(2\pi/d) > 0$ and $\omega_d \notin \mathbb{R}$.  
Proved in \Lean{}.
\end{proof}

The content is that the displacement phase is irreducibly complex: a real-valued character of the
acting group cannot carry it, so the four-valued false branch (a purely real representation of the
WH cocycle) is obstructed, and the projective representation requires genuinely complex cohomology.  
We are deliberately precise about what this is and is not.  
It is the exact non-reality of the phase, which is what blocks the real branch.  
It is \emph{not} the stronger claim that no real fiducial can generate a WH SIC in these dimensions;
that statement is adjacent to the real-equiangular-lines bound and we do not assert it.  
Stating the true obstruction and declining the unproven one is the honest form of the false branch.

\subsection{The dialetheic \texorpdfstring{$B$}{B}-state}

The two branches do not cancel.  
The true branch delivers existence under Stark; the false branch delivers the character obstruction;
and both are theorems.

\begin{theorem}[dialetheic\_fiducial\_paradox]
\label{thm:dialetheia}
Assume the mixed-signature Stark conjecture for $K_{2048}$. Then
\[
\underbracket{\texttt{SICPOVM\_Exists}\,2048}_{\text{true branch}}
\quad\wedge\quad
\underbracket{\neg\,\exists\, r \in \mathbb{R}:\ \omega_{2048} = r}_{\text{false branch}} .
\]
\end{theorem}

Recall the Belnap fixed point $\neg B = B$ of \S\ref{sec:belnap}: the value $B$ absorbs negation
and holds true and false at once.  
Theorem~\ref{thm:dialetheia} places the $d = 2048$ fiducial exactly there.  
The two dimensions the paper closes therefore sit at opposite Belnap values.  
At $d = 12$ the existence node is discharged and the fiducial reads $T$.  
At $d = 2048$ the existence node is open while its obstruction is proved, and the fiducial reads $B$:
not a gap to be apologized for, but the correct four-valued type of an open problem whose skeleton
and whose obstruction are both already theorems.

\subsection{Proved unconditionally (0 sorries)}

The development proves the following results with zero sorries, without any appeal to the axioms that
mark open mathematics.

\begin{itemize}
  \item \textbf{Existence of a SIC-POVM in dimension 12}
    (\texttt{crystal\_forces\_d12\_sic}): exact fiducial constructed in the
    $2048$-dimensional ring $R$, all $143$ overlap identities and the norm
    machine-checked, transferred to $\mathbb{C}^{12}$ along an explicit
    homomorphism.  Depends on no axiom of this development
    (\S\ref{sec:d12}, audit in \S\ref{subsec:audit}).
  \item The Belnap four-valued lattice: $\neg B = B$, no explosion, Frobenius
    closure $\mu \circ \delta = \mathrm{id}$ on $\BFour$.
  \item Product-lattice orbit cardinality: $|\mathcal{O}_{B^{\otimes n}}| = 2^n$
    (\texttt{whOrbit\_card\_eq\_pow2}).
  \item The orbit gap: $2^n < 4^n$ for $n \geq 1$
    (\texttt{product\_lattice\_orbit\_is\_insufficient}).
  \item The group bifurcation: $|\texttt{WHIdx}(n)| = 4^n$
    (\texttt{group\_bifurcation\_lemma}).
  \item The main equivalence: \texttt{hilbert\_embedding\_equiv\_zauner} by \texttt{rfl}.
  \item The Frobenius closure, join-equiangularity, and Born rule in the
    Belnap metric (\texttt{sic\_povm\_belnap\_unconditional}, 22 theorems).
  \item The conditional existence theorem (\texttt{sic\_povm\_exists\_via\_stark}):
    Stark $\Rightarrow$ SIC, from the axioms.
  \item \textbf{The $d = 2048$ transport} (\texttt{ZaunerTransportMap}, 0
    sorries): the unit-norm forward transport
    (\texttt{transport\_preserves\_norm\_sq}), the reduction with
    $\psi\circ\varphi = \mathrm{id}$ (\texttt{reduction\_recovers\_fiducial}),
    the conditional existence (\texttt{zauner\_transport\_theorem}), and the
    character obstruction that $\omega_{2048}$ is non-real
    (\texttt{evalf\_character\_obstruction}), with the dialetheic $B$-state
    conjoining the true and false branches
    (\texttt{dialetheic\_fiducial\_paradox}), Theorems~\ref{thm:transport-unit}--\ref{thm:dialetheia}.
\end{itemize}

\subsection{Axiomatized with known mathematical content}

Table~\ref{tab:axioms} collects the axioms that carry the open mathematical content:
each marks a specific, known gap in the arithmetic geometry rather than a gap in the
formalization strategy.

\begin{table}[h]
\centering
\begin{tabular}{@{}lp{9cm}@{}}
\toprule
\textbf{Axiom} & \textbf{Mathematical content} \\
\midrule
\texttt{zauner\_correspondence} & Galois--Zauner: WH orbit overlaps controlled by the Zauner
  automorphism acting on the Stark unit (Appleby et al.) \\
\texttt{equiangular\_from\_stark\_axiom} & Stark bounds imply equiangularity: the full arithmetic
  geometry of Stark units acting on WH frames \\
\texttt{norm\_of\_normalized\_axiom} & Stark bounds imply the normalized fiducial has the correct
  norm; same arithmetic geometry gap \\
\texttt{stark\_equivalence} & SICPOVM\_Exists $\leftrightarrow$ Stark unit existence
  (Appleby--Flammia--McConnell--Yard structural link; placeholder in formalization) \\
\bottomrule
\end{tabular}
\caption{Honest gap markers: axioms carrying the open mathematical content.}
\label{tab:axioms}
\end{table}

The gap is the translation from the Stark conjecture embedding bounds to the exact equiangularity
and norm equalities.  
This requires the full arithmetic geometry of Galois orbits of Stark units: the machinery of
Shimura varieties and special values of Hecke $L$-functions applied to the specific WH frame setting.  
It is not a gap in the formalization strategy; it is a gap in mathematics.

None of these axioms enters the $d = 12$ existence theorem.  
The audit of \S\ref{subsec:audit} shows that \texttt{crystal\_forces\_d12\_sic} depends only on
\Lean{}'s standard foundations and the compiler trust of \texttt{native\_decide}.  
The table above is the honest ledger of the \emph{general-$d$} chain, not of the $d = 12$ result.

\subsection{Open}

The mixed-signature Stark conjecture for $\Kd$ (Axiom~\ref{ax:stark}) and the Zauner conjecture for
$d = 2^n$, $n \geq 2$ (\texttt{ZaunerConjectureAtPow2}) remain open.  
Theorem~\ref{thm:main-equiv} identifies them as the same problem.  
For $d = 12$ this subsection is empty: every claim about dimension twelve is a machine-checked theorem.  
For $d = 2048$ exactly one item is open: the unconditional existence $\texttt{SICPOVM\_Exists}\,2048$.  
The transport, reduction, conditional reduction to Stark, and character obstruction of \S\ref{sec:d2048}
are all theorems; what is not proved is the single node they surround.

%──────────────────────────────────────────────────────────────────────────────
\section{Discussion}
\label{sec:disc}
%──────────────────────────────────────────────────────────────────────────────

The \texttt{rfl} proof of Theorem~\ref{thm:main-equiv} is the main structural result.  
It states that the Belnap multilattice embedding problem and the Zauner conjecture are definitionally
identical: not logically equivalent (which would require a proof), but the same proposition.

This identification works because both definitions reduce to
\texttt{SICPOVM\_Exists~($2^n$)}, the existence of a WH$(2^n)$-covariant fiducial in
$\mathbb{C}^{2^n}$.  
The Belnap multilattice supplies the structural skeleton (all SIC axioms proved in the Belnap metric,
orbit size $4^n$, Frobenius closure, Born rule); the question of whether that structure can be
realized in $\mathbb{C}^{2^n}$ with the standard WH$(2^n)$ action is exactly what the Zauner
conjecture asks.  
One is the algebraic statement, the other the geometric one; the formalization reveals they are
syntactically the same.

The Stark link (\S\ref{sec:stark}) extends this to number theory.  
The discriminant formula $m_d = (d-3)(d+1)$ is not a coincidence; it is the field whose Galois
theory organizes the WH orbit.  
The Galois--Zauner axiom~\ref{ax:zauner-corr} is the precise statement of the
Appleby--Yadsan-Appleby--Zauner observation \cite{Appleby2013} that the order-3 element of the
Galois group (the ``Zauner automorphism'') controls the SIC overlaps.
Combined with the conditional existence theorem, the SIC-POVM problem becomes, formally, a statement
about special values of Hecke $L$-functions over real quadratic fields.

The present formalization does not prove the Zauner conjecture in generality, nor does it settle
Hilbert's Twelfth Problem.  
It does, however, prove one dimension of Zauner's conjecture outright: $d = 12$, by exact construction,
machine-checked from the ring presentation down to the kernel (\S\ref{sec:d12}).  
Exact $d = 12$ fiducials were known in the literature; what is new is the formal existence proof, to
our knowledge the first for any dimension in any proof assistant, and with it the deletion of the
existence axiom from the development.  
The method deserves a remark of its own: the entire analytic content of the proof was compressed into
a finite set of decidable planks (ring identities under \texttt{native\_decide}, rational sign changes
for the root bracket, divided-difference positivity certificates) plus a single bridge lemma, so the
boundary between computation and analysis is drawn exactly where the kernel can stand on the
computational side.  
A numerical fiducial became an exact ring presentation, and the ring presentation became a theorem;
at no point does the finished proof cite a floating-point quantity.

For the rest, the formalization locates the open content with precision, identifies the three levels
as a single chain, and proves everything in the chain that is not open.  
The remaining axioms are not placeholders for missing proofs; they are the problem.

%──────────────────────────────────────────────────────────────────────────────
\section{Artifact statement}

All modules cited above live in the \texttt{p4ramill} library of the
\texttt{p4rakernel} repository (\url{https://github.com/umpolungfish/p4rakernel}).
Frozen review snapshot: branch \texttt{crystalline/manuscripts3-2026-07-07}, tag
\texttt{crystalline-manuscripts3-v1} (commit \texttt{eea2c0c}); \texttt{main}
continues independently.  This manuscript and its companions
\texttt{witness\_vessel} and \texttt{chrysopoeia\_2048} are frozen in
\texttt{ig-docs} (\url{https://github.com/umpolungfish/ig-docs}) on the same
branch name and tag.

%──────────────────────────────────────────────────────────────────────────────
\section{Acknowledgements}
%──────────────────────────────────────────────────────────────────────────────

The author thanks Harry T.\ Larson for imparting the importance of catching rising problems and
never letting them go \cite{Larson1961}.  
The present formalization is written in that spirit: the Stark--Zauner--Hilbert identification was
not abandoned at the honest gap markers, and dimension $d=12$ was carried from arithmetic prediction
to a machine-checked existence theorem rather than left as a conjectural witness.

%──────────────────────────────────────────────────────────────────────────────
%% Appendix: 1500-digit fiducial witness (auto-generated, do not hand-edit) %%
\appendix
\section{The \texorpdfstring{$d=12$}{d=12} fiducial primary coordinate}
\label{app:witness}

The machine-checked existence proof of \S\ref{sec:d12} produces an explicit fiducial vector in
$\mathbb{C}^{12}$.  
Its primary (zeroth) coordinate is a complex number whose real and imaginary parts are reproduced
below to $1500$ significant decimal digits, exactly as recovered by exact rational computation and
transferred along the ring homomorphism of \S\ref{sec:d12}.  
This is the numerical witness against which the \texttt{native\_decide} audit was run.

\subsection*{Real part}
{\footnotesize\begin{verbatim}
0.1062796056214234892401116553199052855065537600966583358463645722893226773031270658
298768906205486390530515471062546655918150810594007151620739921064457622344183907390
981527630502926525808781095296318291609375737019796692866006233984692833037108933368
557141825426546847294512590606360727608493535003406589845516102474605094572398156211
791281798039542776561629744564194175203915687113788519889443755498620411265436001814
914973032292646202106545011609706046505344547963396007168778128307754312885301096726
162319934008918697128515617656809099978292761800100154277017013499197140932771280829
474710901176233931186576815306983795289327943991170288570563450853952856813572224820
323519766870258180444560783641518280663868403937579668795378352858216503948223011169
697936687708120335054063539968448230213367085847359726613339750534761079437559558177
170458770485092130024050377828296898248742359487812567890426747290383449758739584121
891405313466683715232519962350717633557809679135595933926157129365473633858551371007
959800932525205031235024239569645212695306698728965644437560979625978194688216215727
988558883866587560774612330337201516160788444802262652227620504374101694536195129017
480798390352751711845335405446315804414051104694360997639530534190402662133717527824
453347700164869570420757011587658836186928728835106885477646747032721406405145380391
458059143485376656502679909399982467453578747825974237188668321092959859046845380684
589565644549976727349922297403872284018359670036895362355194841756494931895842349228
42411
\end{verbatim}}

\subsection*{Imaginary part}
{\footnotesize\begin{verbatim}
0.0001630310010432694399587406315542496470085181880359789886468189132905265106516073
242458412714148935119300482376377390198732003194349723579628890157953503274150184941
922432970016102965159540029818814358845133768058607962230856300108004095482775591767
607453478243239371216296563674622826993602940971440231503171851943962519739380612538
628173386359275832075413408173617817250414210122534302281367413780218406525362813767
377242250527245667716771638397748502815187010211458392191084019675813764757830592105
763996740982160026003110762311735045878477445891538375462245767612445866581069506278
477528999281835417095156724276857660199809330614562037038629851055255370186613991140
255317589674458871903863720149657008412573166118787188346680142751797170677947903886
524082866297915466932725341554419267194504939860231321221240228165404147431033338004
194588521299511177701905727538444384985519664781080834434965132700067391423685777638
814268440772310760076171105731976965738049432970937865573746731755532412624301987451
469060456646790760478930383647270817656381589623592381926857668367406314975286153763
506708672554330394678776563753044532036173045952393220064135915855846643399465208690
514984645244764040372464298947848237029742729281882700217734069293190787967268119958
187950201035331429126545813666316372019865413299603495502256362958288261801770979857
300537231820251970793399646987127813352416035330544430067832933876151129206278709923
794881789126474049505755088682279565026575170790581479272730707151904561896233556363
8457817
\end{verbatim}}

%──────────────────────────────────────────────────────────────────────────────
\begin{thebibliography}{99}

\bibitem{Appleby2013}
D.\ M.\ Appleby, H.\ Yadsan-Appleby, and G.\ Zauner,
\textit{Galois automorphisms of a symmetric measurement},
Quantum Inf.\ Comput.\ \textbf{13} (2013), 672--720.

\bibitem{Appleby2017}
M.\ Appleby, S.\ Flammia, G.\ McConnell, and J.\ Yard,
\textit{SICs and algebraic number theory},
Found.\ Phys.\ \textbf{47} (2017), 1042--1059.

\bibitem{AFMY2020}
M.\ Appleby, S.\ Flammia, G.\ McConnell, and J.\ Yard,
\textit{Generating ray class fields of real quadratic fields via complex equiangular lines},
Acta Arith.\ \textbf{192} (2020), 211--233.

\bibitem{Belnap1977}
N.\ D.\ Belnap,
\textit{A Useful Four-Valued Logic},
in: J.\ M.\ Dunn and G.\ Epstein (eds.), \textit{Modern Uses of Multiple-Valued Logic},
Reidel, Dordrecht, 1977, pp.\ 5--37.

\bibitem{Larson1961}
H.\ T.\ Larson,
\textit{Catch a Rising Problem\ldots\ and Never Ever Let it Go},
\emph{IRE Transactions on Engineering Management} \textbf{8}(4) (1961), 173--174.
\url{https://ieeexplore.ieee.org/stamp/stamp.jsp?arnumber=1641382}

\bibitem{Lean4}
L.\ de Moura and S.\ Ullrich,
\textit{The Lean 4 Theorem Prover and Programming Language},
in: CADE-28, Lecture Notes in Comput.\ Sci.\ vol.\ 12699, Springer, 2021, pp.\ 625--635.

\bibitem{p4rakernel}
C.\ L.\ Mills,
\textit{p4rakernel: Lean\,4 paraconsistent kernel fork},
\url{https://github.com/umpolungfish/p4rakernel},
branch \texttt{crystalline/manuscripts3-2026-07-07}, tag
\texttt{crystalline-manuscripts3-v1} (commit \texttt{eea2c0c}), 2026.

\bibitem{Stark1980}
H.\ M.\ Stark,
\textit{$L$-functions at $s=1$, IV: First-order zeroes at $s=1$ for Dirichlet
$L$-functions with real characters},
Advances in Mathematics \textbf{35} (1980), 197--235.

\bibitem{Zauner1999}
G.\ Zauner,
\textit{Quantum designs: Foundations of a noncommutative design theory},
PhD thesis, University of Vienna, 1999.

\end{thebibliography}

\end{document}
